\documentclass[12pt]{article}
\usepackage[T1]{fontenc}
\usepackage[utf8]{inputenc}
\usepackage{lmodern}
\usepackage{textcomp}
\usepackage{enumitem}
\usepackage{geometry}
\usepackage{graphicx}
\usepackage{amsmath, amssymb}
\geometry{margin=1in}
\begin{document}
\begin{center}
Constitutional Law - Graduate Level Assessment 23 \
Total Marks: 40 \
Answer all questions.
\end{center}

\section*{Multiple Choice (10 marks)}
\begin{enumerate}[label=\arabic*.]
\item Neighbor says his home repair contract needs to be discharged because of frustration of purpose. Which of the following is not required for this to be discharged for this reason?
\begin{enumerate}[label=(\alph*)]
    \item A supervening event.
    \item That was reasonably foreseeable at time of contract formation.
    \item Which completely or almost completely destroys the purpose of the contract.
    \item The purpose of the contract was understood by both parties.
\end{enumerate}
\item There is a grand jury proceeding underway for a local businessman. What is not true about grand jury proceedings?
\begin{enumerate}[label=(\alph*)]
    \item The proceedings are conducted in secret.
    \item There is no right to counsel.
    \item There is a right to Miranda warnings.
    \item There is no right to have evidence excluded.
\end{enumerate}
\item There is a voting law that is going to be reviewed by the court. If the law is upheld, what must be true?
\begin{enumerate}[label=(\alph*)]
    \item The law is necessary to achieve a compelling government purpose.
    \item The law is substantially related to an important government purpose.
    \item The law is rationally related to a legitimate government purpose.
    \item The law is substantially related to a legitimate government purpose.
\end{enumerate}
\item Where a client accepts the services of an attorney without an agreement concerning the amount of the fee, there is
\begin{enumerate}[label=(\alph*)]
    \item An implied-in-fact contract.
    \item An implied-in-law contract.
    \item An express contract.
    \item No contract.
\end{enumerate}
\item Which of the following is not a correct form of real evidence?
\begin{enumerate}[label=(\alph*)]
    \item Weapons or implements
    \item marks, scars, wounds
    \item Photographs
    \item Hearsay
\end{enumerate}
\end{enumerate}

\section*{True / False (10 marks)}
\textit{Answer True or False}
\begin{enumerate}[label=\arabic*.]
\setcounter{enumi}{5}
\item True or False: An ex post facto law allows for retroactive punishment based on legislative acts without a trial.
\item True or False: The purpose of the contract must be clearly stated for frustration of purpose to apply.
\item True or False: A defendant charged with first-degree murder is entitled to a list of prospective jurors only upon court order.
\item True or False: The death of a party automatically terminates all contracts involved.
\item True or False: To prove a discriminatory classification, a discriminatory effect alone is insufficient to demonstrate intent to discriminate.
\end{enumerate}

\section*{Long Form Response (20 marks)}
\textit{Answer each question with a clear and complete explanation.}
\begin{enumerate}[label=\arabic*.]
\setcounter{enumi}{10}
\item Discuss the constitutional rationale for Congress enacting a statute that prohibits racial discrimination in real estate transactions, focusing on the enforcement provisions of the Thirteenth Amendment.
\item Discuss the concept of compensatory damages in constitutional law and analyze how they serve to restore the nonbreaching party to the position they would have been in had the contract been fully performed.
\end{enumerate}

\vfill
\noindent\textit{End of Paper}
\end{document}